Because a fixed \(w\in\Delta(\mathcal H)\) exists, the corrected symbol convention implies that \(\mathcal H\) is finite and nonempty. Aggregate the implemented selectors by \(\kappa_{w,a}(O\mid z)=\sum_{v:(v,a)\in O}\kappa_{w,a}(v\mid z)\). Conditional independence of the two auxiliary draws given their common assignment and integration under \(\nu\) give \[R_w(O_0,O_1)=\sum_z\nu(z)\kappa_{w,0}(O_0\mid z)\kappa_{w,1}(O_1\mid z).\tag{1}\] Across replicates these pairs are iid. No factorization across arms is made. Since both target orbits lift the same \(T\), \[P_T(O_0,O_1)=\sum_{v\in\mathcal V}q_{\mathrm{tar}}(v)\mathbf 1\{O_0(v)=O_0,\ O_1(v)=O_1\}.\tag{2}\] Its support is an image of \(\mathcal V\), proving \(L\le\min\{|\mathcal V|,m_0m_1\}\). Fix a preorder pair \(C=(c_0,c_1)\). Aggregate (1) into its ordered class table \(W_{r,s}\) and form the strict prefix table \[H(j,\ell)=\sum_{r<j}\sum_{s<\ell}W_{r,s}.\tag{3}\] A scan of the \(m_0m_1\) cells forms the table. Pushing (2) to class pairs gives \(P_T^C\) and \(L_C\le\min\{L,d_0d_1\}\). For a supported target pair \((j,\ell)\), constant many prefix lookups yield \[\pi^C_{11}=H(j,\ell),\quad \pi^C_{10}=H(j,d_1+1)-H(j,\ell),\quad \pi^C_{01}=H(d_0+1,\ell)-H(j,\ell),\quad \pi^C_{00}=1-\pi^C_{01}-\pi^C_{10}-\pi^C_{11}.\tag{4}\] Conditional on this target class pair, the two rank events are the two Bernoulli-count inequalities with law (4). Thus \(\texttt{lem:four-cell-rank-event-dynamic-program}\) gives \(D_{M,k_0,k_1}(\pi^C(j,\ell))\). Grouping identical exact vectors and averaging with the unchanged target law proves \[\Pr\{\text{both rank events under }C\}=\sum_{\pi\in\Pi_C}\omega_C(\pi)D_{M,k_0,k_1}(\pi),\tag{5}\] and \(U_C\le L_C\). The key \((M,k_0,k_1,\pi_{00},\pi_{01},\pi_{10})\) determines the recurrence value and contains no target weight. Therefore a global balanced tree computes only the \(G\) distinct values, while each occurrence retains its local \(\omega_C\). Local grouping costs \(O(L_C\log(1+L_C))\) comparisons and \(O(L_C)\) arithmetic; global lookup costs \(O((\sum_CU_C)\log(1+G))\). Adding the member scan, class-table scans, target pushforwards, grouping, lookup, and \(G\) recurrences proves the instance-adaptive time display. Replacing \(L_C,U_C\) by their maxima proves the uniform display. The orbit table, coherent target table, rolling recurrence, local exact-key table, and global memo require, respectively, \(O(m_0m_1)\), \(O(L)\), \(O(k_0k_1)\), \(O(U_{\max})\), and \(O(G)\) memory. Finally, the three inherited exact four-cell engines may be selected per key without altering (1)--(5): the zero-cell-safe multinomial formula in \(\texttt{lem:four-cell-multinomial-sum-alternative}\), the corner-oriented rolling recurrence in \(\texttt{lem:corner-oriented-four-cell-rank-evaluator}\), or its positive-base coefficient branch. Put \(\bar k_a=\min\{k_a,M-k_a+1\}\), \[C_{M,k_0,k_1}=\sum_{i=0}^{M-1}\min\{\bar k_0,i+1\}\min\{\bar k_1,i+1\},\] and let \(\mathsf C_+(\pi)\) be the positive-corner coefficient-and-required-marginal charge supplied by \(\texttt{lem:corner-oriented-four-cell-rank-evaluator}\). If \(\mathsf B\) denotes all non-evaluator costs and the distinct exact keys are partitioned into \(\Pi_{\mathrm{MS}}\), \(\Pi_{\mathrm{DP}}\), and \(\Pi_+\) according to the selected branch, sequential exact evaluation costs \[O\!\left(\mathsf B+M+|\Pi_{\mathrm{MS}}|\{M+\mathsf T_{M,\bar k_0,\bar k_1}\}+|\Pi_{\mathrm{DP}}|C_{M,k_0,k_1}+\sum_{\pi\in\Pi_+}\mathsf C_+(\pi)\right).\tag{6}\] The initial \(M\) term safely covers the factorial preprocessing shared by the multinomial keys. If \(u_+(\pi)\) is the coefficient width of the selected positive orientation, the workspace is \[O\!\left(m_0m_1+L+U_{\max}+G+\max\left\{\mathbf 1_{\Pi_{\mathrm{MS}}\ne\varnothing}M,\mathbf 1_{\Pi_{\mathrm{DP}}\ne\varnothing}\bar k_0\bar k_1,\max_{\pi\in\Pi_+}u_+(\pi)\right\}\right),\tag{7}\] with empty-branch terms omitted. Each branch computes the same \(D_{M,k_0,k_1}(\pi)\), and (7) charges the workspace of the branch partition actually selected, not an independently minimized workspace. Since \(C_{M,k_0,k_1}\le Mk_0k_1\), the coarser \(GMk_0k_1\) time and \(k_0k_1\) memory terms in the statement follow by always choosing the rolling branch.